\documentclass[12pt]{article}
\usepackage[english,russian]{babel}

\input aut14.cli
\makeatletter

\usepackage{amssymb}
\usepackage{amsfonts}

\usepackage{setspace}

\usepackage{dsfont}

\usepackage{amsmath,amsthm}

\usepackage{graphics}
\usepackage{graphicx}

\DeclareMathOperator{\LipSh}{LipSh} \DeclareMathOperator{\Int}{Int}
\DeclareMathOperator{\Rep}{Rep} \DeclareMathOperator{\Per}{Per}
\DeclareMathOperator{\Ker}{Ker}
\DeclareMathOperator{\OrientSh}{OrientSh}
\DeclareMathOperator{\dist}{dist} \DeclareMathOperator{\diag}{diag}
\DeclareMathOperator{\StSh}{StSh} \DeclareMathOperator{\Ree}{Re}
\DeclareMathOperator{\Cl}{Cl} \DeclareMathOperator{\Proec}{Pr}
\DeclareMathOperator*{\minn}{min}\DeclareMathOperator*{\KS}{KS}
\DeclareMathOperator{\D}{D}\DeclareMathOperator{\Diff}{Diff}
\DeclareMathOperator{\OrbitSh}{OrbitSh}
\DeclareMathOperator{\PerLin}{PerLin} \DeclareMathOperator{\dd}{d}
\DeclareMathOperator{\Id}{Id}\DeclareMathOperator{\SSeq}{(Seq)}

%\oddsidemargin=21mm

%\textwidth=110mm \textheight=175mm

%\binoppenalty=9999
%\relpenalty=9999
\righthyphenmin=2

%\newcommand{\mybreak}{\smallbreak}
%\newcommand{\Sc}{\large\mdseries\scshape}
%\renewcommand{\thefootnote}{\fnsymbol{footnote})}
%\renewcommand{\paragraph}[1]{\bigbreak{\bfseries #1} }

\date{}
\catcode`@ = 11 \catcode`@ = 12
\def\hiphantom#1{\underline{\phantom{#1}}}

\sloppy \binoppenalty=10000 \relpenalty=10000 \exhyphenpenalty=10000

\begin{document}


%\renewcommand{\normalsize}{\fontsize\@xivpt{14.5}}%\selectfont%
%\abovedisplayskip 11\p@ plus2.5\p@ minus6.5\p@ \belowdisplayskip
%\abovedisplayskip \abovedisplayshortskip  \z@ plus3\p@
%\belowdisplayshortskip  6.5\p@ plus3.5\p@ minus3\p@
%\let\@listi\@listI}

\newcommand{\al}{\mbox{$\alpha$}}
\newcommand{\si}{\mbox{$\sigma$}}
\newcommand{\de}{\mbox{$\delta$}}
\newcommand{\om}{\mbox{$\omega$}}
\newcommand{\De}{\mbox{$\Delta$}}
\newcommand{\ep}{\varepsilon}
\newcommand{\lam}{\mbox{$\lambda$}}
\newcommand{\La}{\mbox{$\Lambda$}}

\newcommand{\RR}{\mbox{$\mathds{R}$}}
\newcommand{\Ss}{\mbox{$\textbf{S}$}}
\newcommand{\TT}{\mbox{$\mathds{T}$}}
\newcommand{\ZZ}{\mbox{$\mathds{Z}$}}
\newcommand{\NN}{\mbox{$\mathds{N}$}}
\newcommand{\QQ}{\mbox{$\mathds{Q}$}}
\newcommand{\Cc}{\mbox{${\bf C}$}}
\newcommand{\Ff}{\mbox{${\cal F}$}}
\newcommand{\Bb}{\mbox{${\cal B}$}}
\newcommand{\Kk}{\mbox{${\cal K}$}}
\newcommand{\Tt}{\mbox{${\cal T}$}}

\newcommand{\KK}{\mbox{${\bf K}$}}

\newcommand{\Ws}{\mbox{$W^s$}}
\newcommand{\Wu}{\mbox{$W^u$}}
\newcommand{\Cone}{\mbox{$\Cc^1$}}


\newcommand{\sref}[1]{(\ref{#1})}
\newcommand{\setg}{\{g(t): \; t \in \RR\}}

\newtheorem{thm}{Теорема }
\newtheorem{lem}{Лемма }
\newtheorem{utv}{Утверждение }
\newtheorem{zam}{Замечание }
\theoremstyle{definition}
\newtheorem{defin}{Определение }
\newtheorem{konstr}{Конструкция }
\newtheorem{primer}{Пример }

\newcommand{\mysection}[1]
{
\medskip
%\begin{flushleft}
\textbf{\large #1} \nopagebreak
%\end{flushleft}
%\smallskip
}

%\newcommand{\mysection}[1]
%{
%\section*{#1}
%}

\sloppy


\makeatletter
\renewcommand{\@oddhead}{\hfill{\large \thepage}\hfill}
\renewcommand{\@oddfoot}{}
\renewcommand{\@evenhead}{\hfill{\large \thepage}\hfill}
\renewcommand{\@evenfoot}{}
\renewcommand{\@makefntext}[1]{\parindent=1em\noindent
 \hbox to 1.8 em {\hss\Large$^{\@thefnmark}$}#1}
\makeatother {\large
\input{%
autoreftitle} \addtocounter{page}{1}
%\input{%
%autoreftitle11}
\input{%
autoreftitle11alt}


\makeatletter
\renewcommand{\@oddhead}{\hfill{\large \thepage}\hfill}
\renewcommand{\@oddfoot}{}
\renewcommand{\@evenhead}{\hfill{\large \thepage}\hfill}
\renewcommand{\@evenfoot}{}
\makeatother

\begin{center}
ОБЩАЯ ХАРАКТЕРИСТИКА РАБОТЫ
\end{center}

\mysection{Актуальность темы.} В диссертации изучается структура множеств диффеоморфизмов гладких многообразий, обладающих так называемы\-ми свойс\-твами отслеживания. 
В настоящее время теория свойств отсле\-жи\-вания является достаточно хорошо разработанным разделом тео\-рии ди\-на\-мических систем. Первые результаты в этом направлении бы\-ли по\-лу\-чены Д.В. Аносовым \cite{An} и Р. Боуэном \cite{Bo}.
 Современное состояние этой теории достаточно полно отражено в монографиях \cite{1,2}. %Данная теория изучает 
%условия, при которых вблизи любой приближенной траектории некоторой динамической системы находится некоторая истинная. 
Грубо говоря, наличие некоторого свойства отслеживания означает, что вблизи любой достаточно точной приближенной траектории дискретной дина\-ми\-чес\-кой системы, порожденной диффеоморфизмом, нахо\-дится ис\-тинная траекто\-рия. В этом случае мы, для прос\-тоты, будем го\-во\-рить, что диффеоморфизм обладает некоторым свойством от\-сле\-живания. Так как термины ''вблизи'' и ''приближенная траектория'' можно формализовывать по-разному, определяют различные свойства отслеживания. Таким образом, если диффеоморфизм обладает не\-ко\-то\-рым свойс\-т\-вом от\-слеживания, то для изучения траекторий точек этого диф\-фео\-мор\-физма можно применять приближенные методы, так как при\-ближенные траектории достаточно адекватно отражают поведение ис\-тин\-ных траекторий. Одним из самых важных и самых первых результатов в теории отслеживания является так называемая shadowing lemma, утверж\-дающая, что определенные свойства отслеживания выполняются в малой окрестности гиперболического множества диф\-фео\-мор\-физ\nobreak ма.

%При изучении различных множеств диффеоморфизмов естественным об\-ра\-зом возникает вопрос о плотности (типичности). Этот воп\-рос пред\-с\-тав\-ляет боль\-шой интерес и для множеств диффео\-мор\-физ\-мов, об\-ла\-дающих некоторым свойством отслеживания. Дей\-с\-т\-ви\-те\-ль\-но, плот\-ность (ти\-пичность) некоторого свойства отслеживания означает, что про\-извольный диффеоморфизм можно так мало возмутить, чтобы он об\-ла\-дал этим свойством отслеживания. 
Боль\-шой интерес представляет вопрос о характеризации множеств диффеоморфизмов, об\-ла\-да\-ющ\-их различными свойствами отслеживания, или внутренностей таких множеств, т.е. о получении необходимых и достаточных условий, при выполнении которых диффеоморфизм (или диффеоморфизм вместе со всеми достаточно близкими к нему диффео\-мор\-физ\-мами) обладает некоторым свойством отслеживания. Кроме того, представляет интерес вопрос о плотности (типичности) множеств диффеоморфизмов, обладающих различными свойствами отслеживания.

\mysection{Цель работы.} Целью данной работы является изучение множеств диффеоморфизмов, обладающих различными свойствами отслеживания, а именно, изучение вопросов о плотности (типичности) этих множеств и о характеризации этих множеств и их внутренностей в терминах теории структурной устойчивости.

\mysection{Методика исследования.} Для получения результатов исполь\-зу\-ют\-ся методы топологии гладких многообразий, дифференциальной гео\-метрии, теории динамических систем~и~др.

\mysection{Научная новизна и значимость работы.} Все результаты дис\-сер\-та\-ции являются новыми. Выделим основные из них:

\begin{itemize}
	\item[--] доказано, что множество $C^1$-диффеоморфизмов глад\-кого замкнутого многообразия, обладающих орбиталь\-ным свойством отслеживания, неплотно в пространстве $C^1$-диффеомор\-физмов с $C^1$-топологией;
	\item[--] доказано, что $C^1$-внутренность множества диффеоморфизмов гладкого замкнутого многообразия, обладающих пе\-ри\-о\-ди\-ческим свойством отслеживания, совпадает с множеством диф\-фео\-морфизмов, обладающих липшицевым периодическим свойс\-т\-вом отслеживания, и с множеством $\Omega$-устойчивых диф\-фео\-мор\-физ\-мов;
	\item[--] доказано, что для любого гомеоморфизма компактного метрического прос\-транства выполняется второе слабое предельное свойство от\-сле\-живания и что множество гомеоморфизмов компактного мет\-рического пространства, обладающих слабым предельным свойством от\-слеживания, обладает следующим свойством: для любого гомео\-морфизма из этого множества неблуждающее множество совпадает с объединением $\omega$-предельных множеств всех траекторий;
	\item[--] доказано, что если гладкое замкнутое многообразие является дву\-мер\-ным, то $C^1$-внутренность множества диффеоморфизмов это\-го многообразия, обладающих слабым предельным свойством отслеживания, совпадает с множеством $\Omega$-устойчивых диф\-фео\-мор\-физмов. 
\end{itemize}

\mysection{Теоретическая и практическая ценность.} Работа носит тео\-ре\-тический характер. Ее результаты важны для теории диффеоморфизмов гладких многообразий и теории ди\-на\-ми\-ческих систем. %Кроме того, ее результаты могут быть использованы при ин\-терпретации результатов численного моделирования траекторий дис\-кретных динамических систем, порожденных диффеоморфизмами.

\mysection{Апробация работы и публикации.} Результаты работы док\-ла\-ды\-ва\-лись на следующих конференциях и семинарах:

\begin{itemize}
	\item[--] Международная конференция ''Современные проблемы математики и механики'', посвященная 70-летию В.А.Садовничего, мех-мат МГУ, март 2009, Москва, Россия;
	\item[--] Topology, Geometry and Dynamics: Rokhlin Memorial, январь 2010, Санкт-Петербург, Россия;
	\item[--] Семинар ''Динамические системы'' (руководители --- академик Д.В.~Аносов и профессор Ю.С.~Ильяшенко), мех-мат МГУ, апрель 2010, Москва, Россия;
	\item[--] Oberseminar Nonlinear Dynamics (руководитель --- профессор B.~Fiedler), Free University Berlin, май 2010, Берлин, Германия;
	\item[--] Петербургский топологический семинар имени В.А. Рохлина (руководитель --- профессор Н.Ю.~Нецветаев), ПОМИ РАН, сентябрь 2010, Санкт-Петербург, Россия;
	\item[--] Семинары по топологии динамических систем (руководитель --- профессор С.Ю.~Пилюгин), кафедра высшей геометрии мат-меха СПбГУ, 2006--2010, Санкт-Петербург, Россия. 
\end{itemize}

Основное содержание диссертации изложено в работах автора \cite{O2}--\cite{O4}. В работе \cite{O4} автору диссертации принадлежит доказательство теоремы о совпадении $C^1$-внутренности множества диффеоморфизмов, обладающих периодическим свойством отслеживания, и множества $\Omega$-устойчивых диф\-фео\-морфизмов, а также леммы о том, что если для диффеоморфизма вы\-пол\-няется липшицево периодическое свойство отслеживания, то лю\-бая периодическая точка гиперболична. %Все остальное принадлежит соав\-то\-рам. 
Статьи \cite{O2,O4} опубликованы в изданиях, включенных в перечень ВАК на момент публикации.

\mysection{Структура и объем работы.} Диссертационная работа объемом 132 машинописных страницы состоит из введения (первой главы), трех глав основной части и списка литературы, содержащего 34 наименования. Вторая глава включает 7 параграфов, третья и четвертая главы --- по 3 параграфа. Седьмой параграф второй главы состоит из четырех пунктов.

\medskip
\begin{center}
СОДЕРЖАНИЕ РАБОТЫ
\end{center}
\medskip

Во введении (первой главе) даны определения исследуемых в дис\-сер\-тации свойств отслеживания, сформулированы основные результаты дис\-сер\-тации, и приведено сравнение этих результатов с соответствующими результатами других математиков. 

Во второй главе изучается проблема плотности (типичности) свойств отслеживания.

Основным объектом данной работы являются диффеоморфиз\-мы глад\-ких замкнутых многообразий. Однако для определения свойств от\-сле\-жи\-вания гладкость не нужна, поэтому различные свойства отслеживания будут определяться для гомеоморфизмов метрических пространств. 

В данной работе под типичным понимается свойство, выполняющееся для всех отображений из некоторого множества II категории по Бэру, а под плотным --- выполняющееся для всех отображений из некоторого плотного множества. В работе рассматриваются пространство гомеомор\-физ\-мов $H(M)$ ком\-пактного метрического пространства $M$ с $C^0$-топологи\-ей %ин\-ду\-цируемой метрикой
%$$d_{C^0}(f,g)=\max_{p\in M}\max\left(\mbox{dist}(f(p),g(p)),\mbox{dist}(f^{-1}(p),g^{-1}(p))\right).$$
%Пусть $M$ --- это гладкое замкнутое многообразие с римановой метрикой dist. Будем обозначать через $Df(x)$ дифференциал диффеоморфизма $f$ многообразия $M$ в точке $x$. Как обычно, будем обозначать через $T_xM$ ка\-сательное пространство в точке $x$ многообразия $M$. Будем обозначать че\-рез $|\cdot|$ норму в пространстве $T_xM$, порожденную метрикой dist.
%Для сокращения изложения, будем считать, что многообразие $M$ вложено в евклидово пространство $\mathbb{R}^{\mathcal{N}}$ достаточно большой размерности. В этом случае касательное пространство естественным образом отождествляется с линейным подпространством пространства $\mathbb{R}^{\mathcal{N}}$.
%Будем обозначать че\-рез $\mbox{Diff}^1(M)$ прос\-транство диффеоморфизмов гладкого замкнутого мно\-го\-об\-ра\-зия $M$ с $C^1$-топологией, индуцируемой метрикой
%$$d_{C^1}(f,g)=\max_{p\in M}\left(\mbox{dist}(f(p),g(p)),\max_{v\in T_pM, |v|=1}|Df(p)v-Dg(p)v|\right).$$
и пространство $\mbox{Diff}^1(M)$ $C^1$-диффеоморфизмов гладкого замкнутого многообразия $M$ с $C^1$-топологией.
% Отметим, что пространства $H(M)$ и $\mbox{Diff}^1(M)$ являются простанствами Бэра, поэтому если некоторое свойство отображений является типичным относительно одной из этих топологий, то оно является и плотным от\-но\-сительно этой топо\-логии.

Пусть $f$ --- гомеоморфизм метрического пространства $(M,\mbox{dist})$.  
Траекторией точки $p$ гомеоморфизма $f$ называется множество
$$O(p,f)=\left\{f^k(p)\mid k\in\mathbb{Z}\right\}.$$
%Аналогично вводятся понятия положительной и отрицательной полу\-траекторий гомеоморфизма $f$:
%$$O_{+}(p,f)=\left\{f^k(p)\mid k\in\mathbb{Z}, k\geq 0\right\},$$
%$$O_{-}(p,f)=\left\{f^k(p)\mid k\in\mathbb{Z}, k\leq 0\right\}.$$ 

%\textbf{Определение 1.} 
Будем называть последователь\-ность $\xi=\{x_k\}$ 
\textit{$d$--псевдотраек\-торией} гомеоморфизма $f$, если 
$$\mbox{dist}(x_{k+1},f(x_k))<d\qquad\mbox{при}\ k\in\mathbb{Z}.$$

%%Таким образом, $d$--псевдотраектория является одной из возможных формализациий понятия приближенной траектории. 

%\textbf{Определение 2.} 
%Будем говорить, что для гомеоморфизма $f$ выполняется \textit{свойство} POTP \textit{(свойство отслеживания псевдотраекторий, pseudo orbit tracing property)}, если для любого положительного $\epsilon$ существует такое положительное $d$, что для любой $d$-псевдотраектории $\xi=\{x_k\}$ найдется такая точка $q\in M$, что
%\begin{equation}
%\label{1}
%\mbox{dist}(x_k,f^k(q))< \epsilon\quad\mbox{ при }k\in\mathbb{Z}.
%\end{equation}
%Иными словами, свойство POTP состоит в том, 
%что всякая ''достаточно точная'' псевдотраектория ''поточечно близка'' к некоторой истинной траектории. В работе мы будем обозначать одним и тем же символом как некоторое свойство динамических систем, так и множество всех систем, обладающих этим свойством.
%Мы будем говорить, что точка $q$ $\epsilon$-отслеживает псевдотраекторию $\xi$, если выполняются неравенства $(\ref{1})$.

Будем обозначать через $N(\epsilon,A)$ $\epsilon$-окрестность множества $A\subset M$. 
В работе \cite{3} вводится определение орбитального свойства отслеживания (OSP, orbital shadowing property), являющегося основным объектом ис\-сле\-до\-ва\-ния второй главы. %и слабого свойства отслеживания (WSP, weak shadowing property).

%\textbf{Определение 3.} 
Будем говорить, что для гомеоморфизма $f$  выполняется \textit{орбитальное свойство отслеживания}, если для любого положительного $\epsilon$ существует такое положительное $d$, что для любой $d$--псевдотраек\-тории $\xi$ найдет\-ся та\-кая точка $q\in M$, что
%\begin{equation}
%\label{3}
$$\xi\subset N(\epsilon,O(q,f))\qquad\mbox{и}\qquad O(q,f)\subset N(\epsilon,\xi).$$
%\end{equation}

%\textbf{Определение 4.} 
%Будем говорить, что для гомеоморфизма $f$ выполняется \textit{слабое свойство отслеживания}, если для любого положительного $\epsilon$ существует такое положительное $d$, что для любой $d$-псевдотраектории $\xi$ найдется такая точка $q\in M$, что
%$$\xi\subset N(\epsilon,O(q,f)).$$

%Свойство OSP является ослабленным аналогом свойства POTP --- вместо того чтобы требовать близость ''в каждый момент времени'' точки псевдотраектории $x_k$ и точки истинной траектории $f^k(q)$, требуется, 
%чтобы были близки множества точек 
%псевдотраектории $\xi=\{x_k\}$ и траектории $O(q,f)$. Слабое свойство отслеживания WSP является ослабленным вариантом свойства OSP --- требуется лишь, чтобы множество точек ''достаточно точной'' псевдотраектории $\xi$ содержалось в малой окрестности некоторой истинной траектории $O(q,f)$. 
%Мы будем говорить, что траектория точки $q$ орбитально $\epsilon$-отслеживает псевдотраекторию $\xi$, если выполняются неравенства~$(\ref{3})$.

%Иными словами, наличие орбитального свойства отслеживания оз\-на\-ча\-ет, что любая достаточно точная псевдотраектория и некоторая истинная траектория близки как множества, т.е. в смысле расстояния по Хаусдорфу.

%С.Ю. Пилюгин и О.Б. Пламеневская (см. $\cite{4}$) доказали типичность свойства отслеживания псевдотраекторий в пространстве $H(M)$, если пространство $M$ является гладким замкнутым многообразием. Отметим что, из $C^0$-типичности свойства отслеживания псевдотраекторий следует $C^0$-типичность орбитального и слабого свойств отслеживания. Ч. Бонатти, Л.Дж. Диац и Дж.~Туркат \cite{5} показали, что свойство отслеживания псевдотраекторий является неплотным относительно $C^1$-топологии в пространстве $\mbox{Diff}^1(M)$, а С.~Кровизье \cite{6} установил, что слабое свойство отслеживания $C^1$-плотно. %(см. также работу С.Ю. Пилюгина, К. Сакая и О.А. Тараканова \cite{PT}). 

В диссертации перечисляются известные результаты о плот\-ности и типичности различных свойств отслеживания. Не перечисляя их все, упомянем только, что из работы С.Ю. Пи\-люгина и О.Б. Пламеневской \cite{4} следует типичность орбитального свойс\-тва отслеживания в пространстве $H(M)$.
Основным результатом второй главы является следующая теорема:

\textbf{Теорема 1.} \textit{Существует такая область} $W\subset\mbox{Diff}^1(S^2\times S^1)$\textit{, что любой диффеоморфизм $f\in W$ не обладает свойством} OSP\textit{.}  

%Обозначим через $\Sigma^2$ метрическое пространство всех двухсторонних последовательностей $0$ и $1$, а через $\sigma$ --- сдвиг Бернулли. В статье \cite{17} вводятся следующие определения.

%Зафиксируем два диффеоморфизма $f_0$ и $f_1$ окружности $S^1$. Бу\-дем на\-зывать \textit{ступенчатым косым произведением} такое отображение $G:\Sigma^2\times S^1\mapsto\Sigma^2\times S^1$, что
%$$G(\omega,\phi)=(\sigma(\omega), f_{\omega_0}(\phi))\qquad\mbox{при}\ \omega\in\Sigma^2,\phi\in S^1,$$
%где $\omega_0$ --- символ, стоящий у последовательности $\omega$ на нулевой позиции.

%\textbf{Определение 13.} 
%Зафиксируем произвольный набор диффеоморфизмов $f_{\omega}$ ок\-руж\-нос\-ти~$S^1$, параметризованный двухсторонними пос\-ле\-до\-ватель\-нос\-тями $\omega\in\Sigma^2$. Будем называть \textit{мягким косым произведением} такое отоб\-ражение $G:\Sigma^2\times S^1\mapsto\Sigma^2\times S^1$, что
%$$G(\omega,\phi)=(\sigma(\omega), f_{\omega}(\phi))\qquad\mbox{при}\ \omega\in\Sigma^2,\phi\in S^1.$$

В параграфе 2.1 описана схема доказательства теоремы~1. Для доказательства теоремы 1 используется идея, восходящая к А.С. Го\-ро\-дец\-кому и Ю.С. Ильяшенко: строить пример в клас\-се час\-тич\-но-ги\-пер\-болических мягких косых произведений (см. \cite{17}). %Для этого рас\-сма\-т\-ри\-ва\-ет\-ся некоторое ступенчатое косое произведение $G_0$ над сдви\-гом Бернулли со слоем --- окружностью. Реализовав сдвиг Бер\-нул\-ли как отображение подковы Смейла, достаточно быстро сжимаю\-щее и рас\-тягивающее по сравнению с динамикой на слое, мы уви\-дим, что это\-му ступенчатому косому произведению соответствует глад\-кое отоб\-ражение с локально-максимальным частич\-но-гипербо\-ли\-чес\-ким мно\-жес\-твом, центральные слои в котором -- окружности. Более того, как следует из работ М.В. Хирша, К.К. Пью, М. Шуба и А.С. Городецкого (см.~\cite{12,13}), при 
%$C^1$-малых возмущениях такой гладкой реализации частич\-но-гиперболическое множество сохраняется, по-прежнему оставаясь произведением окружности на подкову 
%Смейла. Так как индивидуально глад\-кие цен\-тральные слои гельдерово зависят от точки на ба\-зе (т.е. на подкове Смейла), сужение возмущенного отображения на это час\-тич\-но ги\-пер\-болическое множество является гельдеровым мяг\-ким ко\-сым произведением. 
В качестве искомой области $W$ из теоремы 1 берется достаточно малая $C^1$-окрестность некоторого ступенчатого косого произведения $G_0$. Эта окрестность берется, в частности, настолько малой, что любому диффеоморфизму из области $W$ соответствует некоторое мягкое косое произведение. Далее, доказывается, что теорема 1 сводится к теореме $1^{\prime}$ о свойствах мягких косых произведений. Эта теорема утверждает, что для гельдеровых мягких косых произве\-де\-ний, ''дос\-та\-точ\-но близких'' к косому произведению $G_0$, не выполняется орбитальное свойство отслеживания.

%\textbf{Теорема $\mbox{1}^{\prime}$.} \textit{Для гельдеровых мягких косых произве\-де\-ний, ''дос\-та\-точ\-но близких'' к косому произведению $G_0$, не выполняется}~OSP\textit{.}

Доказательство теоремы $1^{\prime}$ разбивается на два случая: слу\-чай (A1) и случай (A2). Случай (A1) возникает, если существуют гиперболическая периодическая точка с одномерным неустойчивым многообразием и гиперболическая периодическая точка с одномерным устойчивым многообразием, у которых эти многообразия пересекаются. %Здесь и далее под ги\-пер\-бо\-ли\-чес\-ки\-ми периодическими точками мягких косых произведений по\-ни\-ма\-ют\-ся периодические точки, которым соответствуют гиперболические пе\-риодические точки диффеоморфизмов, а под устойчивыми (не\-ус\-той\-чи\-вы\-ми) многообразиями --- множества, соответствующие ус\-той\-чи\-вым (не\-ус\-тойчивым) многообразиям этих гиперболичес\-ких пе\-ри\-о\-ди\-чес\-ких точек. Под\-чер\-к\-нем, что, несмотря на название, ус\-той\-чи\-вые и не\-ус\-той\-чи\-вые мно\-гообразия гиперболических периодических точек мягких косых произведений многообразиями не яв\-ля\-ют\-ся. 
Да\-лее, с помощью основной леммы строятся псевдотраектории, которые не могут отслеживаться никакой точной траекторией.
Слу\-чай (A2) возникает при отсутствии таких пе\-ре\-сечений. 
В этом случае строится такая псевдотраектория, что лю\-бая ор\-би\-та\-ль\-но отслеживающая ее точная 
траектория должна оказаться ге\-те\-ро\-кли\-нической, соединяющей две гиперболические периодичес\-кие точ\-ки с соответственно одномерным неустойчивым и одномерным ус\-той\-чивым многообразиями. Предположение о наличии орбитального свой\-с\-тва отслеживания противоречит условию случая (A2).
%В этом слу\-чае мы построим такую псевдотраекторию, что лю\-бая ор\-би\-та\-ль\-но отслеживающая ее точная 
%траектория должна оказаться ге\-те\-ро\-кли\-нической, соединяющей две гиперболические периодичес\-кие точ\-ки с соответственно одномерным неустойчивым и одномерным ус\-той\-чивым многообразиями. Предположение о наличии орбитального свой\-с\-тва отслеживания будет противоречить сделанному предположению об отсутствии пересечений. Подчеркнем, что случай (A2) является основным, и именно в нем в полную силу используется техника мягких косых про\-изведений.

В параграфе 2.2 описываются основные свойства косых произведений и обсуждается сведение теоремы 1 к теореме $\mbox{1}^{\prime}$. В параграфе~2.3 фор\-мулируется лемма 1 (основная лемма), содержащая дос\-та\-точ\-ное условие на псевдотраекторию, при котором любая орбитально от\-сле\-живающая ее траектория лежит на устойчивом (или 
неустойчивом) мно\-гообразии гиперболической периодической точки. 
В параграфе 2.4 по\-казано, что доказательство теоремы сводится к двум случаям и разбирается случай (A1).
В параграфе 2.5 доказаны вспомогательные лем\-мы о свойствах сдвига Бернулли, которые утверждают, что если не\-ко\-торая последовательность содержит не слишком мало нулей под\-ряд, то ее траектория не попадает в определенные области. В пара\-г\-ра\-фе~2.6 разбирается случай (A2) с точностью до леммы~6, которая поз\-во\-ля\-ет построить псевдотраекторию с необходимыми свойст\-вами. До\-казательству леммы 6 посвящен параграф~2.7.%, который состоит из четырех пунктов. Отметим, что в доказательстве используется техника из работы~\cite{17}. 

В третьей главе изучаются периодические свойства отслеживания и их связь со структурной устойчивостью. 

Пусть $f$ --- гомеоморфизм метрического пространства $(M,\mbox{dist})$.  
%Можно рассмотреть аналог свойства отслеживания псев\-до\-тра\-ек\-то\-рий для периодических траекторий и пе\-рио\-ди\-чес\-ких псев\-до\-траекторий. 
В работе \cite{K} дается определение периодического свой\-ства отслеживания.

%\textbf{Определение 4.} 
Будем говорить, что для гомеоморфизма $f$ выполняется \textit{пе\-ри\-оди\-чес\-кое свойство отслеживания} PerSh \textit{(periodic shadowing property)}, если для любого положительного $\epsilon$ существует такое положительное~$d$, что для любой периодической $d$-псев\-дотраектории $\xi$ найдется такая периодическая точка $q$, что выполняется соотношение %$(\ref{1})$.
$$\mbox{dist}(x_k,f^k(q))<\epsilon\quad\mbox{при }k\in\mathbb{Z}.$$
%В определении свойства POTP можно накладывать условия на зависимость числа $\epsilon$ от $d$. Это приводит к %следующему определению. 
%
%\textbf{Определение 5.} Будем говорить, что для гомеоморфизма $f$ выполняется \textit{свойство} LipSP %\textit{(Lipschitz shadowing property, липшицево свойство отслеживания)}, если существуют такие положительные числа %$L$ и $d_0$, что для любой $d$-псевдотраектории $\xi=\{x_k\}$ с $d\leq d_0$ найдется такая точка $q\in M$, что
%\begin{equation}
%\label{2}
%\mbox{dist}(x_k,f^k(q))\leq Ld\quad\mbox{ при }k\in\mathbb{Z}.
%\end{equation}

%Наряду с периодическим свойством отслеживания можно рас\-смат\-ривать и его липшицев аналог.

%\textbf{Определение 6.} 
Будем говорить, что для гомеоморфизма $f$ выполняется \textit{липшицево пе\-ри\-оди\-ческое свойство отслеживания} LipPerSh\textit{,} если существуют та\-кие положительные числа $L$ и $d_0$, что для любой периодической $d$-псев\-до\-тра\-ек\-тории с $d\leq d_0$ найдется такая периодическая точка $q$, что выполняется со\-отношение %$(\ref{2})$.
$$\mbox{dist}(x_k,f^k(q))\leq Ld\quad\mbox{ при }k\in\mathbb{Z}.$$
 
При изучении свойств отслеживания диффеоморфизмов гладких зам\-к\-ну\-тых многообразий одним из важных подходов является переход к $C^1$-внут\-рен\-ностям. Будем обозначать через $\mbox{Int}^1(P)$ внутренность некоторого под\-мно\-жества $P$ множества диффеоморфизмов гладкого зам\-кнутого мно\-го\-образия $M$ относительно $C^1$-топологии. Как обыч\-но, будем обоз\-на\-чать через $\Omega S$ множество $\Omega$-устойчивых диффео\-мор\-физмов. В дис\-сер\-та\-ции перечисляются известные результаты о ха\-рак\-теризации различных свойств отслеживания в терминах тео\-рии струк\-тур\-ной устойчивости. 
%В работе \cite{10} К. Сакай доказал, что $C^1$-внутренность множества диффеоморфизмов, обладающих свойством отслеживания псевдотраекторий, совпадает с множеством структурно устойчивых диффеоморфизмов. 
В третьей главе доказывается, что $C^1$-внут\-рен\-ность мно\-жес\-т\-ва диффеоморфизмов, обладающих периодическим свойст\-вом от\-сле\-жи\-вания, совпадает с множеством $\Omega$-устойчивых диф\-фео\-мор\-физ\-мов.
%Существуют не структурно устойчивые диффеоморфизмы, обладающие свойством отслеживания псевдотраекторий. 
%(см., например, \cite{P}). 
От\-ме\-тим, что существуют не $\Omega$-ус\-той\-чивые диффеоморфизмы, об\-ла\-да\-ющие периодическим свойством отслеживания.  %С.Ю. Пилюгин и С.Б. Тихомиров в работе \cite{11} доказали, что структурная устойчивость и липшицево свойство отслеживания эквивалентны. %Также отметим, что С.Ю. Пилюгин доказал (см. \cite{P}) эквивалентность структурной устойчивости и так называемого вариационного свойства отслеживания.
В третьей гла\-ве также до\-ка\-зы\-вается эквивалентность свойства $\Omega$-устойчивости и липшицева периоди\-чес\-ко\-го свойства отслеживания. Таким образом, основной результат третьей главы можно сформулировать~так:

\textbf{Теорема 2.} \textit{Если $M$ --- гладкое замкнутое многообразие, то} 
$\mbox{Int}^1(\mbox{PerSP})=\mbox{LipPerSP}=\Omega S$\textit{.}

В параграфе 3.1 доказывается, что множество $\Omega$-ус\-той\-чи\-вых диффеоморфизмов является подмножеством мно\-жес\-т\-ва диф\-фео\-мор\-физмов, обладающих липшицевым периодическим свойством отслеживания. В параграфе 3.2 доказывается, что $C^1$-внутрен\-ность мно\-жества диффеоморфизмов, обладающих периодическим свойством от\-слеживания, содержится в множестве $\Omega$-устойчивых диффеоморфизмов. Наконец, в параграфе 3.3 доказывается, что множество диффеоморфизмов, обладающих липшицевым периодическим свойством отслежива\-ния, яв\-ля\-ет\-ся подмножеством множества $\Omega$-устойчивых диффеоморфизмов. Для этого сначала доказывается гиперболичность всех периодических точек, потом их равномерная гиперболичность, потом плотность периодических точек в неблуждающем множестве и, наконец, выполнение условия от\-сут\-с\-т\-вия циклов. 

В четвертой главе изучаются  слабые предельные свойства от\-сле\-жи\-ва\-ния. В параграфе 4.1 даются определения предельных свойств от\-сле\-жи\-ва\-ния и рассматриваются их простейшие свойства.

%Пусть $M$ --- метрическое пространство с метрикой dist, $f:M\mapsto M$ --- гомеоморфизм. 
Пусть $f$ --- гомеоморфизм метрического пространства $(M,\mbox{dist})$.
Определим следующее свойство последовательнос\-ти $\xi=\{x_k\}_{k>0}$:
$$\mbox{dist}(x_{k+1},f(x_k))\longrightarrow0\quad\mbox{при }k\rightarrow+\infty.\eqno(*)$$
%Последовательность, обладающая свойством $(*)$, является приближенной траекторией, которая становится ''все точнее'' на бесконечности.
%В работе \cite{7} введено следующее определение.
%
%\textbf{Определение 7.} 
%Будем говорить, что для гомеоморфизма $f$ выполняется \textit{свойство} LmSP \textit{(предельное свойство отслеживания, limit shadowing property)}, если для любой последовательности, для которой выполняется свойство $(*)$, найдется такая точка $q\in M$, что
%$$\mbox{dist}(x_k,f^k(q))\longrightarrow0\quad\mbox{при }k\rightarrow+\infty.$$ 
%
%Наличие предельного свойства отслеживания означает, что если приближенная траектория ''становится все точнее'' на бесконечности, то к ней ''стремится'' некоторая точная траектория. По аналогии с обычным свойством отслеживания предельное свойство отслеживания можно ослабить, потребовав не ''сходимости'' приближенной траектории к точной, а совпадения множеств предельных точек. Обозначим через  $\omega(\xi)$ множество всех предельных точек последовательности $\xi=\{x_k\}_{k>0}$.  
%В работе \cite{8} вводится следующее определение.  
%
%\textbf{Определение 8.} 
%Будем говорить, что для гомеоморфизма $f$ выполняется \textit{орбитальное предельное свойство отслеживания} OLmSP, если для любой последовательности $\xi\!=\!\{x_k\}$, для которой выполняется свойство $(*)$, найдется такая точка $q\in\! M$, что
% $$\omega(\xi)=\omega(q).$$
% 
%Нетрудно видеть, что $$\mbox{LmSP}\subset\mbox{OLmSP}.$$
%
%В работе \cite{8} вводится следующее ослабление орбитального предельного свойства отслеживания.
В работе \cite{8} вводится следующее свойство.

%\textbf{Определение 9.}
 Будем говорить, что для гомеоморфизма $f$ выполняется \textit{слабое пре\-де\-ль\-ное свойс\-тво отслеживания} WLmSP, если для любой после\-до\-ва\-те\-льности $\xi\!=\!\{x_k\}_{k\geq0}$, для которой выполняется свойство $(*)$, найдется такая точ\-ка $q\in\! M$, что 
%\begin{equation}
%\label{0.3}
$$\omega(\xi)\subset\omega(q),$$
%\end{equation}
где через $\omega(\xi)$ обозначено множество всех предельных точек последовательности $\xi$.

%В работе \cite{3} наряду со свойством WSP определяется и следующее свойство. 
%
%\textbf{Определение 10.}
% Будем говорить, что для гомеоморфизма $f$ выполняется \textit{второе слабое свойство отслеживания} 2WSP, если для любого положительного $\varepsilon$ найдется такое положительное $d$, что для любой $d$-псевдотраектории $\xi$ найдется такая точка $q\in\! M$, что $$O(q,f)\subset N(\varepsilon,\xi).$$
%
%Нетрудно видеть, что свойство 2WSP является ослаблением свойства OSP.
%По аналогии с обычными свойствами отслеживания, исходя из определения орбитального предельного свойства отслеживания, мы вводим еще одно слабое предельное свойство отслеживания.
%Можно определить еще одно слабое предельное свойство отсле\-жи\-вания.
 
% \textbf{Определение 11.}
  Будем говорить, что для гомеоморфизма $f$ выполняется \textit{вто\-рое сла\-бое предельное свойс\-тво отслеживания} 2WLmSP, если для любой пос\-ле\-до\-вательности $\xi\!=\!\{x_k\}$, для которой выполняется свойство $(*)$, найдется такая точка $q\in\! M$, что 
%\begin{equation}
%\label{0.1}
$$\omega(q)\subset\omega(\xi).$$
%\end{equation}

Можно ввести аналоги сформулированных предельных свойств от\-сле\-жи\-ва\-ния, в которых $k\rightarrow+\infty$ заменено на $k\rightarrow-\infty$ (в главе 4 приводятся эти определения).

%В работе \cite{3} доказано, что любой гомеоморфизм $f$ компактного метрического пространства $M$ обладает свойством 2WSP. 
%В параграфе 4.1 мы доказываем аналог этого утверждения для свойства 2WLmSP ---
%доказано, что любой гомеоморфизм $f$ компактного метрического пространства $M$ обладает свойством 2WLmSP. 
В параграфе 4.1 доказывается, что любой гомеоморфизм $f$ ком\-пак\-т\-ного мет\-рического пространства $M$ обладает свойством 2WLmSP.

 При изучении динамических систем особый интерес вызыва\-ет струк\-тура неблуждающего множества. Как обычно, обозначим через $\Omega(f)$ мно\-жес\-тво всех неблуждающих точек гомеоморфизма~$f$.
Хорошо известно, что для гомеоморфизма $f$
$$\Omega(f)\supset \bigcup_{p\in M}\omega(p).$$
В параграфе 4.2 приводится пример гомеоморфизма $f$ для которого множества $\Omega(f)$ и $\bigcup_{p\in M}\omega(p)$ не совпадают. Кроме того, там доказывается, что если для гомеоморфизма $f$ выполняется свойство WLmSP, то
$$\Omega(f)= \bigcup_{p\in M}\omega(p).$$

В диссертации перечисляются известные результаты о ха\-рак\-те\-ризации $C^1$-внутренностей множеств диффеоморфизмов, обладающих предельными свойствами отслеживания, в тер\-ми\-нах теории структурной устойчивости. %Для краткости мы их здесь не приводим.
% С.Ю. Пилюгин в работе \cite{8} доказал, что 
% $$\mbox{Int}^1(\mbox{OLmSP})=\mbox{Int}^1(\mbox{LmSP})=\Omega S.$$ 
% К. Сакай в работе \cite{S} показал, что если $\dim M=2$, то $\mbox{Int}^1(\mbox{WSP})\subset\allowbreak\subset\Omega S$. Кроме того, К. Сакай отметил, что его результат не обобщается на случай многообразий большей размерности. 
Основным результатом четвертой главы является следующая тео\-рема, доказательство которой приведено в параграфе 4.3.

\textbf{Теорема 3.} \textit{Если $M$ --- гладкое замкнутое многообразие и $\dim M=2$, то} 
$$\mbox{Int}^1(\mbox{WLmSP})=\Omega S.$$

В четвертой главе приводится пример, показывающий, что теорема~3 не обобщается на слу\-чай многообразий более высокой размерности.
Для ''отрицатель\-ных'' пре\-дельных свойств отслеживания выполняются аналогичные утвер\-ж\-де\-ния. Идея доказательсва состоит в использовании леммы Хаяши и Аоки (см.~\cite{H,A}), которая утверждает, что $C^1$-внутренность множества диф\-фео\-морфизмов, все периодические точки которых гиперболичны, совпадает с множеством $\Omega$-устойчивых диффеоморфизмов. Далее, мы показываем, что диффеоморфизм, имеющий негиперболическую периодическую точку, можно так $C^1$-мало возмутить, чтобы не выполнялось слабое предельное свойство отслеживания.
%при определенных условиях наличие негиперболической периодической точки противоречит наличию слабого предельного свойства отслеживания.

\makeatletter
%Стиль библиографии
\renewcommand{\@biblabel}[1]{#1.\hfill}

\renewcommand{\refname}{Список цитируемой литературы}

\begin{thebibliography}{99}
\bibitem{An}\textit{Аносов Д. В.} Об одном классе инвариантных множеств гладких
динамических систем // Труды 5-й Межд. конф. по нелин. колеб. Киев. 1970. Т. 2. С. 39--45.
%\bibitem{12} \textit{Городецкий А. С.} Минимальные аттракторы и частично ги\-пер\-бо\-ли\-чес\-кие инвариантные множества динамических систем. Текст диссертации, МГУ им. Ломоносова, мех-мат, 2001, 77 с.
\bibitem{17}\textit{Городецкий А. С., Ильяшенко Ю. С.} Некоторые свойства косых про\-изведений над подковой и соленоидом. // Труды мат. инст. им. В. А. Стеклова РАН. 2000. Т. 231. С. 96--118.
\bibitem{A}\textit{Aoki N.}, The set of Axiom A diffeomorphisms with no cycle, Bol.
Soc. Brasil. Mat. (N.S.), Vol. 23, 1992, P. 21--65.
%\bibitem{5} \textit{Bonatti Ch., Diaz L. J., Turcat G.} Pas de ''shadowing lemma'' pour les dynamiques partiellement hyperboliques // C. R. Acad. Sci. Paris, Ser I, Math. 2000. Vol. 330. P. 587--592.
\bibitem{Bo}\textit{Bowen R.} Equilibrium States and the Ergodic Theory of Anosov
Diffeomorphisms. Lect. Notes. Math. Vol. 470. Springer-Verlag.
Berlin. 1975. 108~p.
%\bibitem{6} \textit{Crovisier S.} Periodic orbits and chain-transitive sets of $C^1$-diffeomorphisms // Publ. Math. de LI'H$\acute{\mbox{E}}$S. 2006. Vol. 104. \No 1. P. 87--141.
%\bibitem{7} \textit{Eirola T., Nevanlinna O., Pilyugin S. Yu.} Limit shadowing property // Numer. Funct. Anal. Optim. 1997. Vol. 18. P. 75--92.
%\bibitem{14}\textit{Gorodetski A.} Regularity of central leaves of partially hyperbolic sets and its applications // Izv. RAN (The Newslet. of the Rus. Acad. of Science). 2006. Vol. 70. \No 6. P. 52--78.
\bibitem{H}\textit{Hayashi S.} Diffeomorphisms in $\mathcal{F}^1(M)$ satisfy Axiom A // Ergod. Theory
Dyn. Syst., Vol. 12, 1992, P. 233--253.
%\bibitem{13}\textit{Hirsch M. W., Pugh C. C., Shub M.} Invariant Manifolds, Lect. Notes in Math., Vol. 583. Springer-Verlag. Berlin-New York. 1977. 149 p.
\bibitem{K}\textit{Ko$\acute{\mbox{s}}$cielniak P.}, On genericity of shadowing and periodic shadowing
property // J. Math. Anal. Appl., Vol. 310, 2005, P. 188--196.
\bibitem{2} \textit{Palmer K.} Shadowing in Dynamical Systems. Methods and Applications. Dordrecht-Boston-London. 2000. 299 p.  
\bibitem{8} \textit{Pilyugin S. Yu.} Sets of diffeomorphisms with various limit shadowing
properties // J. Dyn. Differ. Equat. 2007. Vol. 19. P. 747--775.
\bibitem{1} \textit{Pilyugin S. Yu.} Shadowing in Dynamical Systems, Lect. Notes in Math. Vol. 1706. Springer-Verlag. Berlin. 1999. 283 p. 
%\bibitem{P}\textit{Pilyugin S. Yu.} Variational shadowing // Discr. Contin. Dyn. Syst.
%(accepted).
\bibitem{4} \textit{Pilyugin S. Yu., Plamenevskaya O. B.} Shadowing is generic // Topol. and its Appl. 1999. Vol. 97. P. 253--266.
\bibitem{3} \textit{Pilyugin S. Yu., Rodionova A. A., Sakai K.} Orbital and weak shadowing properties // Discr. Contin. Dyn. Systems. 2003. Vol.~9. P.~287--308.
%\bibitem{PT}\textit{Pilyugin S. Yu., Sakai K., Tarakanov O. A.} Transversality properties and
%$C^1$-open sets of diffeomorphisms with weak shadowing // Discr. Contin. Dyn. Systems, 2006, Vol. 16, N 4, P. 871--882.
%\bibitem{11} \textit{Pilyugin S. Yu., Tikhomirov S. B.} Lipschitz shadowing implies
%structural stability (to appear).
%\bibitem{10}\textit{Sakai K.} Pseudo orbit tracing property and strong
%transversality of diffeomorphisms of closed manifolds // Osaka J.
%Math. 1994. Vol. 31. P. 373--386.
%\bibitem{S}\textit{Sakai K.} Diffeomorphisms with weak shadowing // Fund. Math., 2001. Vol.~168, P. 53--75.
\end{thebibliography}

\renewcommand{\refname}{Публикации автора по теме диссертации}

\begin{thebibliography}{9}
\bibitem{O2}\textit{Осипов А. В.} Неплотность орбитального свойства отслеживания относительно $C^1$-топологии // Алгебра и анализ, 2010, Т.~22, С. 127--163.
\bibitem{O1}\textit{Осипов А. В.} Слабые предельные свойства отслеживания динами\-ческих систем на двумерных многообразиях // Электр. журн. Диф. ур-я и проц. упр-я., 2006. N\ 4, С. 1--16.  
\bibitem{O3}\textit{Осипов А. В.} $C^1$-неплотность орбитального свойства отслеживания // Зап. межд. конф. совр. пробл. матем. и прил., Москва, Унив. книга, 2009, С. 183--184.
\bibitem{O5}\textit{Osipov A. V.} Lipschitz Periodic Shadowing // Topology, Geometry and Dynamics: Rokhlin Memorial, St. Petersburg, 2010, P. 58--59.
\bibitem{O4}\textit{Osipov A. V., Pilyugin S. Yu., Tikhomirov S. B.} Periodic shadowing and $\Omega$-stability // Reg. Chaotic Dynamics, 2010, Vol. 15, N\ 2--3, P.406--419.
\end{thebibliography}

\newpage

\makeatletter
\renewcommand{\@oddhead}{\hfill{\large }\hfill}
\renewcommand{\@oddfoot}{}
\renewcommand{\@evenhead}{\hfill{\large }\hfill}
\renewcommand{\@evenfoot}{}
\makeatother

\clearpage \vspace*{\fill}

%\hrule
%\begin{center}
%\small Подписано к печати 26.12.08 Формат $60 \times 84$ $1/16$.\\
%Бумага офсетная. Гарнитура Times. Печать цифровая. Печ. л. 0,7.\\
%Тираж 100 экз. Заказ 4367.
%\end{center}
%\hrule
%\begin{center}
%\small Отпечатано в отделе оперативной полиграфии химического
%факультета СПбГУ\\
%198504, Санкт-Петербург, Старый Петергоф, Университетский пр., 26\\
%Тел.: (812)428-4043, 428-6919
%\end{center}
%\vspace{2cm}
%\bigskip
%\bigskip
%\bigskip
%\bigskip
%\bigskip
%\bigskip

\hrule
\begin{center}
\small
Подписано в печать \dots \dots. Формат $60 \times 81$ 1/16.\\
Бумага офсетная. Печать офсетная. \\
Усл. печ. листов 0,7. Тираж 100 экз. Заказ N \dots \\
\end{center}
\hrule
\begin{center}
\small
ЦОП типографии издательства СПбГУ\\
199061, С-Петербург, Средний пр., д. 41.
\end{center}
\hrule

\end{document}
